\section*{Hoofdstuk \ref{ch:g12:contdist} --- Continue toevalsvariabelen}

\begin{solution}{exo:g12:contdist:1}
$\P(X \leq 5) = \frac{5}{15} = \frac13$;
$\P(4 \leq X \leq 10) = \frac{6}{15} = \frac25$;
$\E(X) = \frac{15}{2} = 7.5$ min;
$\sigma(X) = \frac{15}{\sqrt{12}} = \frac{5\sqrt3}{2} \approx 4.3$ min.
\end{solution}

\begin{solution}{exo:g12:contdist:2}
$f \geq 0$ en
$\int_0^2 \frac38 t^2\,\dd t = \frac38\left[\frac{t^3}{3}\right]_0^2
= \frac38 \times \frac83 = 1$: een dichtheid.
\[
\E(X) = \int_0^2 \frac38 t^3\,\dd t = \frac38 \times \frac{16}{4} = \frac32,
\qquad
\P(X \geq 1) = \int_1^2 \frac38 t^2\,\dd t = \frac{8 - 1}{8} = \frac78 .
\]
\end{solution}

\begin{solution}{exo:g12:contdist:3}
\emph{1.} $\E(X) = \frac{1}{0.2} = 5$ jaar;
$\P(X > 5) = \eu^{-0.2\times5} = \eu^{-1} \approx 0.37$.

\emph{2.} Wegens de geheugenloosheid (\cref{thm:g12:contdist:memoryless}) is
$\pcond{X > 3}{X > 8} = \P(X > 5) = \eu^{-1} \approx 0.37$: de drie jaar
dienst veranderen niets.
\end{solution}

\begin{solution}{exo:g12:contdist:4}
$\P(X > m) = \eu^{-\lambda m} = \frac12$ geeft
$m = \frac{\ln 2}{\lambda} \approx \frac{0.69}{\lambda}$, kleiner dan
$\E(X) = \frac1\lambda$. De exponentiële dichtheid heeft een lange rechterstaart:
enkele ongewoon lange levensduren trekken het \emph{\omterm{def:g11:stat:mean}{gemiddelde}} boven de
\emph{\omterm{def:g11:stat:median}{mediaan}}, de waarde die de helft van de populatie overtreft. (Dat is de
halveringstijd uit \cref{ex:g12:exp:halflife}: de helft van de atomen overleeft
ze, hoewel de \omterm{def:g11:stat:mean}{gemiddelde} levensduur langer is.)
\end{solution}

\begin{solution}{exo:g12:contdist:5}
$168 = \mu - \sigma$ en $182 = \mu + \sigma$: ongeveer $68\%$.
$189 = \mu + 2\sigma$: daarboven ligt de helft van de resterende
$100 - 95.4 = 4.6\%$, dus ongeveer $2.3\%$.
$154 = \mu - 3\sigma$: ongeveer $\frac{100 - 99.7}{2} = 0.15\%$.
\end{solution}

\begin{solution}{exo:g12:contdist:6}
$\P(X < 500) \leq 2.5\%$ betekent dat $500$ minstens twee
\omterm{def:g11:stat:variance}{standaardafwijkingen} onder het \omterm{def:g11:stat:mean}{gemiddelde} moet liggen (de \omterm{def:g12:contdist:normal}{normale verdeling}
laat $2.3\% \approx 2.5\%$ onder $\mu - 2\sigma$; met de gebruikelijke $1.96$
verloopt de redenering identiek):
\[
\mu - 2\sigma \geq 500 \iff \mu \geq 500 + 8 = 508 \text{ g}.
\]
De machine moet gemiddeld op $508$ g ingesteld worden --- de prijs van de
waarborg is $8$ g ``gratis'' product per zak.
\end{solution}

\begin{solution}{exo:g12:contdist:7}
\emph{1.} Met $p = \frac16$ en $n = 1200$ is
$\sqrt{\frac{p(1-p)}{n}} = \sqrt{\frac{5/36}{1200}} \approx 0.0108$, dus is
het schommelingsinterval op $95\%$
\[
\frac16 \pm 1.96 \times 0.0108 \approx \intcc{0.146}{0.188}.
\]

\emph{2.} De waargenomen \omterm{def:g10:stats:series}{frequentie} $\frac{260}{1200} \approx 0.217$ ligt er
ver buiten: de zuiverheidshypothese wordt op het niveau $5\%$ verworpen. De
afwijking bedraagt ongeveer $4.6$ \omterm{def:g11:stat:variance}{standaardafwijkingen} --- voor een normale
\omterm{def:g10:numbers:approx}{benadering} een \omterm{def:g10:proba:distribution}{kans} van de orde $10^{-6}$, veel doorslaggevender dan de grens
$\leq 4.6\%$ van Bienaymé--Chebyshev (\cref{exo:g12:sums:7}).
\end{solution}

\begin{solution}{exo:g12:contdist:8}
\emph{1.} De marge $\frac{1}{\sqrt n} \leq 0.02$ vraagt $n \geq 2500$ mensen.

\emph{2.} Om scores die $1$ punt uit elkaar liggen te onderscheiden, moet de
marge ruim onder $0.5$ punt blijven, wat volgens de vereenvoudigde formule
$n \geq \left(\frac{1}{0.005}\right)^2 = 40\,000$ vraagt --- voor de meeste
peilingen al onhaalbaar. Erger nog: de statistische marge verrekent alleen de
fout van het \emph{bemonsteren}; systematische vertekeningen (niet
representatieve steekproeven, weigeringen, verschuivingen op het laatste
ogenblik) krimpen niet naarmate $n$ groeit en overtreffen doorgaans één punt.
Een race met $1$ punt \omterm{def:g11:seq:arithmetic}{verschil} is werkelijk te nipt om te beslechten.
\end{solution}

\begin{solution}{exo:g12:contdist:9}
\emph{1.} Omdat $\alpha > 0$, geldt $c \leq \alpha X + \beta \leq d
\iff \frac{c - \beta}{\alpha} \leq X \leq \frac{d - \beta}{\alpha}$, dus
\[
\P(c \leq Y \leq d)
= \int_{(c-\beta)/\alpha}^{(d-\beta)/\alpha} f(t)\,\dd t .
\]
De substitutie $y = \alpha t + \beta$ (dat wil zeggen: de oppervlakte in de
veranderlijke $y$ lezen, met $t = \frac{y - \beta}{\alpha}$ en
$\dd t = \frac{\dd y}{\alpha}$) maakt hiervan
$\int_c^d \frac{1}{\alpha} f\left(\frac{y-\beta}{\alpha}\right)\dd y$: de
\omterm{def:g11:func:function}{functie} $g(y) = \frac1\alpha f\left(\frac{y-\beta}{\alpha}\right)$ is een
dichtheid van $Y$.

\emph{2.} Met $f = \varphi$, $\alpha = \sigma$ en $\beta = \mu$:
\[
g(y) = \frac{1}{\sigma}\,\varphi\!\left(\frac{y - \mu}{\sigma}\right)
= \frac{1}{\sigma\sqrt{2\pi}}\,
\exp\left(-\frac{(y-\mu)^2}{2\sigma^2}\right),
\]
en dat is dus de dichtheid van
$\mathcal N(\mu, \sigma^2) = \mu + \sigma\,\mathcal N(0,1)$.
\end{solution}

\begin{solution}{pb:g12:contdist:1}
\textbf{1.} $\P(\text{wachttijd} \leq 10) = \frac{10}{30} = \frac13$;
$\E = 15$ min; $\sigma = \frac{30}{\sqrt{12}} \approx 8.7$ min.

\textbf{2.} $\int_0^1 2x\,\dd x = 1$: een dichtheid.
$\E(X) = \int_0^1 2x^2\,\dd x = \frac23$;
$\P\left(X > \frac12\right) = \int_{1/2}^1 2x\,\dd x = 1 - \frac14 =
\frac34$.

\textbf{3.} $\P(T > 15) = \eu^{-15/10} \approx 0.22$. De \omterm{def:g11:stat:median}{mediaan}:
$\eu^{-m/10} = \frac12$, dus $m = 10\ln 2 \approx 6.9$ min --- minder dan het
\omterm{def:g11:stat:mean}{gemiddelde} $10$, omdat de lange rechterstaart van de \omterm{def:g12:contdist:expo}{exponentiële verdeling}
(zeldzame reusachtige wachttijden) het \omterm{def:g11:stat:mean}{gemiddelde} omhoogtrekt terwijl de
\omterm{def:g11:stat:median}{mediaan} bij de ``typische'' wachttijd blijft.

\textbf{4.} Wegens de geheugenloosheid is
$\P(T > 35 \mid T > 20) = \P(T > 15) \approx 0.22$: de twintig reeds
uitgezeten minuten kopen niets --- de stroom veroudert niet.

\textbf{5.} (a) exponentieel (het radioactief verval is het model waarvoor de
geheugenloosheid uitgevonden werd); (b) uniform op $\intcc{0}{8}$ (een
uurrooster met een willekeurige verschuiving); (c) exponentieel; (d) geen van
beide: een versleten lamp brandt in het volgende uur \emph{waarschijnlijker}
door dan een nieuwe, dus $\P(T > s + t \mid T > s) < \P(T > t)$: veroudering
is in strijd met de stelling over de geheugenloosheid.

\textbf{6.} $68\,\%$, $95\,\%$, $99.7\,\%$: de drie ijkintervallen.

\textbf{7.} $z = \frac{190 - 172}{8} = 2.25$: bovenstaart $\approx 1.2\,\%$.

\textbf{8.} $z = 2$: ongeveer $2.3\,\%$ --- ruwweg één persoon op $44$.

\textbf{9.} De specificatie ligt op $\pm 2\sigma$: ongeveer $95.4\,\%$ komt
erdoor en $4.6\,\%$ wordt afgekeurd --- op grote schaal rampzalig. Op
$\pm 6\sigma$ zakt het uitvalpercentage tot ongeveer twee op een
\emph{miljard}: zes sigma koopt het recht om massaal te produceren zonder
massaal te keuren.

\textbf{10.} $\sigma = 5$; met continuïteitscorrectie is
$\P \approx \P\left(\abs Z \leq \frac{5.5}{5}\right) = \P(\abs Z \leq 1.1)
\approx 0.73$. De stelling maakt de bord van Galton uit
\cref{pb:g11:binom:1} glad: de trap van vakjes wordt de continue klok.

\textbf{11.} $n \approx \frac{1.96^2 \times 0.25}{0.005^2} \approx 38\,400$
mensen --- en bij die omvang overheersen de fouten die \emph{niet} van het
bemonsteren komen (het steekproefkader, weigeringen, leugens) de marge: een
race met één punt \omterm{def:g11:seq:arithmetic}{verschil} ligt buiten het bereik van eerlijk peilen, en
daarom zeggen ernstige bureaus dan ``te nipt om te beslechten''.

\textbf{12.} Geheugenloosheid: vanaf het ogenblik dat je aankomt, is de
resterende wachttijd exponentieel met \omterm{def:g11:stat:mean}{gemiddelde} $10$ --- de volle $10$
minuten, niet $5$. De ``\omterm{def:g11:stat:mean}{gemiddelde} tussentijd $10$'' van het uurrooster
misleidt: jouw wachttijd heeft dezelfde verdeling als een hele tussentijd.

\textbf{13.} Willekeurige aankomsten vallen evenredig met de lengte van de
tussentijd, dus met \omterm{def:g10:proba:distribution}{kans} $\frac{15}{20} = \frac34$ in een lange. Verwachte
wachttijd: $\frac34 \times \frac{15}{2} + \frac14 \times \frac52 =
5.625 + 0.625 = 6.25$ min --- meer dan de naïeve $5$, hoewel de \omterm{def:g11:stat:mean}{gemiddelde}
tussentijd $10$ is. De schuldige: het \emph{naar lengte vertekende
bemonsteren} --- lange tussentijden vangen meer reizigers.

\textbf{14.} De wachttijd achterwaarts en die voorwaarts zijn elk exponentieel
met \omterm{def:g11:stat:mean}{gemiddelde} $10$ (de geheugenloosheid werkt vanaf je aankomst in beide
richtingen): de tussentijd waarin je zit, duurt gemiddeld $20$ minuten --- het
dubbele van de typische tussentijd. De inspectieparadox in één zin: \emph{het
\omterm{def:g10:numbers:interval}{interval} dat je toevallig inspecteert, is geen typisch \omterm{def:g10:numbers:interval}{interval}, want je viel
met grotere \omterm{def:g10:proba:distribution}{kans} in een groot \omterm{def:g10:numbers:interval}{interval}.}

\textbf{15.} \omterm{def:g11:stat:mean}{Gemiddelde} klas: $\frac{9 \times 10 + 110}{10} = 20$ studenten.
Het \omterm{def:g11:stat:mean}{gemiddelde} dat de studenten beleven:
$\frac{90 \times 10 + 110 \times 110}{200} = \frac{13\,000}{200} = 65$
studenten. De brochure drukt $20$ af; drie op vijf studenten zitten in de klas
van $110$ en beleven de $65$.

\textbf{16.} Populaire mensen staan op vele vriendenlijsten, dus ``een vriend''
bemonsteren is naar lengte vertekend in het voordeel van de gezelligen ---
jouw vrienden worden getrokken uit de lange tussentijden van het sociale
uurrooster.

\textbf{17.} Voor $t \geq 0$:
$\P(T > t) = \P(-10\ln U > t) = \P\left(U < \eu^{-t/10}\right) =
\eu^{-t/10}$: precies de overlevingsfunctie van de \omterm{def:g12:contdist:expo}{exponentiële verdeling} ---
één \omterm{def:g12:exp:ln}{logaritme} zet de uniforme ruis van de computer om in eender welke
wachttijd.

\textbf{18.} Elke uniforme variabele heeft \omterm{def:g11:stat:mean}{gemiddelde} $\frac12$ en \omterm{def:g12:randvar:exp}{variantie}
$\frac{1}{12}$: de som van twaalf heeft \omterm{def:g11:stat:mean}{gemiddelde} $6$ en \omterm{def:g12:randvar:exp}{variantie} $1$ ---
dus levert het recept \omterm{def:g11:stat:mean}{gemiddelde} $0$ en \omterm{def:g12:randvar:exp}{variantie} $1$ op. Twaalf
onafhankelijke kleine duwtjes opgeteld: het mechanisme van Galton, en met de
zegen van De Moivre--Laplace is het histogram voor het oog al klokvormig.

\textbf{19.} Over het model: markten zijn geen sommen van vele
\emph{onafhankelijke} kleine oorzaken --- paniek correleert alles (de les van
2008 uit de vorige opgave) en brengt dikke staarten voort die geen enkele
normale dichtheid bezit. Wanneer de gegevens $10^{-137}$ fluisteren, antwoordt
de catechismus van de statisticus: \emph{het model wordt verworpen, niet de
dag}.

\textbf{20.} Uniform: alle posities gelijk, het eerlijke ``ik ken alleen de
grenzen''. Exponentieel: het risico dat nooit veroudert --- verval,
aankomsten, klikken. Normaal: de democratische klok van vele kleine
onafhankelijke oorzaken --- lengtes, fouten, \omterm{def:g11:stat:mean}{gemiddelden}. De angel: wat je
bemonstert, is vertekend door de manier waarop je erop botste --- bussen,
klassen, vrienden. En de boog: het kind dat pakjes sap telde, heeft tien
volumes vol opgaven later de menigte met een kromme gemeten; de epsilon, de
\omterm{def:g12:integ:area}{integraal} en de centrale limietstelling wachten in de universitaire volumes om
uit te leggen \emph{waarom} de klok voor alles luidt.
\end{solution}
